\documentclass{article}
\usepackage[margin=1in]{geometry}
\usepackage{longtable}
\begin{document}

\textbf{Legacy note:} This split-by-person draft complements the canonical roadmap in \texttt{Documentation/SRS.tex}.\\

\section*{Phase 6 Plan - Person 2 (Person B)}

\textbf{Source of truth:} \texttt{Documentation/SRS.tex}, Section 11.5, ``Phase 6: Train and Deploy ML v1''.\\
\textbf{Interpretation note:} Person 2 owns the modeling and evaluation logic for Phase 6. The job is to turn replayable episode data into a first learned ranker, compare it to required baselines, calibrate the outputs, and define usable alert cutoffs.\\
\textbf{Implementation note:} Person 2 should treat the model as one candidate ranking policy among several until baseline comparisons and calibration are convincing. Shadow deployment does not count as success by itself.

\textbf{Phase duration:} 2-4 weeks\\
\textbf{Phase goal:} Train, compare, calibrate, and document ML v1.

\section*{Owned Deliverables (From SRS)}
\begin{enumerate}
\item Build the episode training dataset.
\item Train the first LightGBM or CatBoost ranker.
\item Compare against the required baselines.
\item Calibrate outputs by market liquidity bucket and category.
\item Define the cutoffs for \texttt{WATCH}, \texttt{ACTIONABLE}, and \texttt{CRITICAL}.
\item Produce a model card describing failure modes and data limitations.
\end{enumerate}

\section*{Locked Implementation Decisions}
\begin{itemize}
\item The training dataset must be derived from replayable episode artifacts and stable feature versions.
\item Baseline comparisons are mandatory; the model must outperform wallet-unaware or simpler ranking baselines to justify adoption.
\item Calibration must be segmented by liquidity bucket and category, not assumed globally uniform.
\item Alert cutoffs must be tied to calibrated score behavior and operational tolerance, not arbitrary score bands.
\item The model card must document failure modes, data limits, and known blind spots honestly.
\end{itemize}

\section*{Concrete Execution Breakdown}
\begin{longtable}{p{0.28\linewidth}p{0.32\linewidth}p{0.30\linewidth}}
\textbf{Work Item} & \textbf{Artifact} & \textbf{Done Condition} \\
Training dataset & Episode-level ML dataset & Replay-derived rows are stable, labeled, and versioned enough for repeatable experiments \\
First ranker & LightGBM or CatBoost model artifact & One first-pass model can score episodes or alerts consistently against held-out data \\
Baseline comparison & Baseline-comparison report & Model performance is contextualized against required non-ML baselines \\
Calibration and cutoffs & Calibration report plus alert thresholds & Score outputs are interpretable and mapped to operational severity buckets \\
Model card & Failure-mode and limitation memo & Risks, blind spots, and data constraints are documented honestly \\
\end{longtable}

\section*{Execution Sequence (No Day Timeline)}
\textbf{Block 1: Training dataset construction}
\begin{itemize}
\item Build the episode training dataset from replay-derived features and stable labels.
\item Define train/validation/test splits that respect the Phase 5 holdout logic.
\item Record dataset hashes and feature-schema versions for every training run.
\end{itemize}

\textbf{Block 2: First ranker and baseline comparisons}
\begin{itemize}
\item Train the first LightGBM or CatBoost ranker.
\item Compare it against the required baselines, especially wallet-unaware or simpler ranking policies.
\item Reject any result that looks better only because of leakage or inconsistent feature handling.
\end{itemize}

\textbf{Block 3: Calibration and alert cutoffs}
\begin{itemize}
\item Calibrate outputs by market liquidity bucket and category.
\item Define the cutoffs for \texttt{WATCH}, \texttt{ACTIONABLE}, and \texttt{CRITICAL}.
\item Validate that the cutoffs make operational sense in shadow-mode alert review.
\end{itemize}

\textbf{Block 4: Model card and handoff}
\begin{itemize}
\item Produce a model card describing failure modes, limitations, and rollback concerns.
\item Hand the calibrated artifact and threshold recommendations to Person 1 for registry and shadow deployment.
\item Summarize whether ML v1 deserves continued shadowing or should be revised first.
\end{itemize}

\section*{Handoffs to Person 1}
\begin{itemize}
\item Trained model artifact and exact dataset hash.
\item Baseline-comparison and calibration reports.
\item Recommended alert thresholds for \texttt{WATCH}, \texttt{ACTIONABLE}, and \texttt{CRITICAL}.
\item Model card with failure modes and deployment caveats.
\item Any feature-contract issues discovered during training.
\end{itemize}

\section*{Gate 6 Checklist (Person 2 Ownership)}
\begin{itemize}
\item Episode training dataset is stable and versioned.
\item First ranker is trained and compared against required baselines.
\item Calibration exists by liquidity bucket and category.
\item Score cutoffs are defined for operational severities.
\item Model card documents failure modes and data limitations.
\end{itemize}

\end{document}
